% F18：本书原创矢量示意；依据所在节的定义、证明或明确模型。
\begin{figure}[htbp]
\centering
\begin{tikzpicture}[x=1cm,y=1cm,>=stealth,every node/.style={font=\small},line width=.65pt]

\draw[->,gray](0,0)--(6,0)node[right]{\(x\)};
\fill[AnalysisBlue!10](2,0)rectangle(4,2);
\draw[AnalysisBlue,thick](.5,0)..controls(1.2,0)and(1.2,2)..(2,2)--(4,2)..controls(4.8,2)and(4.8,0)..(5.5,0);
\draw[very thick,orange!80!black](2.3,0)--(3.7,0);
\node at(3,-.4){集合及其开邻域};\node at(3,3.25){容许函数在邻域上至少为一};
\node[align=left]at(9.5,2.1){零阶能量：\(r^d A\)\\[8pt]梯度能量：\(r^{d-p}B\)};
\node at(9.5,.3){\(E_p(\eta_r)=r^d A+r^{d-p}B\)};
\node at(9.5,-.45){相对容量只取梯度项};

\end{tikzpicture}
\caption[容量的容许函数与过渡能量]{容量的容许函数与过渡能量。平台与过渡区对应容许性和梯度代价。曲线只示意截断形状，容量仍是对所有容许函数及开邻域取下确界所得的量。}
\label{b1:fig:visual-f18}
\end{figure}
